% !TEX program = pdflatex
\documentclass[11pt]{article}

% ==========================
% Encoding / language (pdflatex)
% ==========================
\usepackage[utf8]{inputenc}
\usepackage[T1]{fontenc}
\usepackage[english]{babel}
\usepackage{lmodern}
\usepackage{microtype}

% ==========================
% Page layout
% ==========================
\usepackage[margin=1in]{geometry}

% ==========================
% Math
% ==========================
\usepackage{amsmath,amssymb,amsthm,mathtools}

% ==========================
% Lists
% ==========================
\usepackage{enumitem}

% ==========================
% Hyperref
% ==========================
\usepackage[hidelinks]{hyperref}

% ==========================
% Macros
% ==========================
\newcommand{\Pb}{\mathbb{P}}
\newcommand{\E}{\mathbb{E}}
\newcommand{\R}{\mathbb{R}}
\newcommand{\F}{\mathcal{F}}
\newcommand{\1}{\mathbf{1}}
\newcommand{\dd}{\,\mathrm{d}}
\newcommand{\norm}[1]{\left\lVert #1\right\rVert}

% ==========================
% Theorem-like environments
% ==========================
\theoremstyle{definition}
\newtheorem{exercise}{Exercise}[section]
\newtheorem{remark}{Remark}[section]

% ==========================
% Exercise utilities
% ==========================
% Optional short header at the beginning of each lecture section
\newcommand{\lecturesummary}[1]{%
  \begin{quote}\small\textbf{Focus.} #1\end{quote}
}

% A compact enumerate style often useful inside exercises
\newlist{exparts}{enumerate}{1}
\setlist[exparts,1]{label=\textbf{(\alph*)}, itemsep=0.25em, topsep=0.4em, leftmargin=2em}

% ==========================
% Document metadata
% ==========================
\title{Stochastic Calculus --- Training Problem Set (Lectures 1--6)}
\author{}
\date{}

\begin{document}
\maketitle
\vspace{-1.2em}

\noindent\textbf{Instructions.}
\begin{itemize}[leftmargin=1.6em]
  \item This sheet is meant for practice before the midterm/final.
\end{itemize}

% ==========================================================
\section{Lecture 1 --- Option Pricing : Motivations and First Mathematical and Numerical tools}

% --- Put ~10--15 exercises here ---
% =====================================================
%  Exercise 1
% =====================================================
\begin{exercise}[Quadratic transformation of a normal variable]
Let $X\sim \mathcal N(0,1)$ and define $Y:=X^2$.
\begin{enumerate}[label=\textbf{(\alph*)}, itemsep=0.25em]
  \item Compute the cdf $F_Y(y)=\mathbb P(Y\le y)$ for $y\ge 0$.
  \item Deduce the density $f_Y$.
  \item Identify the distribution of $Y$ (name and parameters).
  \item Compute $\mathbb E[Y]$ and $\mathrm{Var}(Y)$.
\end{enumerate}
\end{exercise}

% =====================================================
%  Exercise 2
% =====================================================
\begin{exercise}[Characteristic function: inversion and identification]
Let $X$ be a real-valued random variable with characteristic function
$$
\varphi_X(t)=\mathbb E\!\left[e^{itX}\right]=\frac{1}{1+t^2},\qquad t\in\mathbb R.
$$
\begin{enumerate}[label=\textbf{(\alph*)}, itemsep=0.35em]

  \item \textbf{Warm-up checks.}
  \begin{enumerate}[label=\textbf{(\roman*)}, itemsep=0.2em]
    \item Verify that $\varphi_X(0)=1$ and that $\varphi_X$ is even and real-valued.
    \item Show that $|\varphi_X(t)|\le 1$ for all $t\in\mathbb R$.
    \item Compute $\displaystyle \lim_{|t|\to\infty}\varphi_X(t)$.
  \end{enumerate}

  \item \textbf{Integrability.}
  Show that $\varphi_X\in L^1(\mathbb R)$ by computing
  $$
  \int_{-\infty}^{\infty}\left|\varphi_X(t)\right|\,dt
  =
  \int_{-\infty}^{\infty}\frac{1}{1+t^2}\,dt.
  $$
  (You may recall $\int_{-\infty}^{\infty}\frac{1}{1+t^2}\,dt=\pi$.)

  \item \textbf{Recovering a density (Fourier inversion).}
  Assume $\varphi_X\in L^1(\mathbb R)$. Then $X$ admits a continuous density $f_X$ given by
  $$
  f_X(x)=\frac{1}{2\pi}\int_{-\infty}^{\infty} e^{-itx}\,\varphi_X(t)\,dt.
  $$
  \begin{enumerate}[label=\textbf{(\roman*)}, itemsep=0.2em]
    \item Using the identity
    $$
    \int_{-\infty}^{\infty} e^{-itx}\,\frac{1}{1+t^2}\,dt
    = \pi e^{-|x|},
    $$
    deduce that
    $$
    f_X(x)=\frac12 e^{-|x|},\qquad x\in\mathbb R.
    $$
    \item Check that $f_X$ integrates to $1$.
  \end{enumerate}

  \item \textbf{Identification of the law.}
  \begin{enumerate}[label=\textbf{(\roman*)}, itemsep=0.2em]
    \item Identify the distribution with density $f_X(x)=\frac12 e^{-|x|}$ (name and parameters).
    \item Compute $\mathbb P(X\ge 0)$ and explain why it must be $1/2$.
  \end{enumerate}

  \item \textbf{Moments: existence and computation.}
  \begin{enumerate}[label=\textbf{(\roman*)}, itemsep=0.2em]
    \item Does $\mathbb E[X]$ exist? Justify carefully.
    \item Compute $\mathbb E[X]$.
    \item Compute $\mathbb E[|X|]$ and $\mathrm{Var}(X)$.
  \end{enumerate}

  \item \textbf{Smoothness of $\varphi_X$ and moments (optional).}
  \begin{enumerate}[label=\textbf{(\roman*)}, itemsep=0.2em]
    \item Compute $\varphi_X'(t)$ and check whether $\varphi_X'(0)$ exists.
    Recall: if $\mathbb E[|X|]<\infty$, then $\varphi_X'(0)=i\,\mathbb E[X]$.
    \item Compute $\varphi_X''(0)$ and discuss what it suggests about $\mathbb E[X^2]$.
    (Recall: if $\mathbb E[X^2]<\infty$, then $\varphi_X''(0)=-\mathbb E[X^2]$.)
  \end{enumerate}

  \item \textbf{Comparison with the Gaussian case (optional).}
  Let $Z\sim\mathcal N(0,\sigma^2)$ with characteristic function
  $$
  \varphi_Z(t)=\exp\!\left(-\frac{\sigma^2 t^2}{2}\right).
  $$
  \begin{enumerate}[label=\textbf{(\roman*)}, itemsep=0.2em]
    \item Compare the decay of $\varphi_X(t)=\frac{1}{1+t^2}$ and $\varphi_Z(t)$ as $|t|\to\infty$.
    \item Explain qualitatively what this implies about the tails of $X$ versus $Z$.
    \item Which one has heavier tails? Give one quantitative argument (e.g.\ tail asymptotics).
  \end{enumerate}

\end{enumerate}
\end{exercise}

% =====================================================
%  Exercise 3
% =====================================================
\begin{exercise}[Gaussian conditioning (explicit formulas)]
Let $(X,Y)$ be jointly Gaussian with
$$
\mathbb E[X]=\mu_X,\quad \mathbb E[Y]=\mu_Y,\quad
\mathrm{Var}(X)=\sigma_X^2,\quad \mathrm{Var}(Y)=\sigma_Y^2,\quad
\mathrm{Cov}(X,Y)=\rho\sigma_X\sigma_Y.
$$
\begin{enumerate}[label=\textbf{(\alph*)}, itemsep=0.25em]
  \item Show that
  $$
  \mathbb E[X\mid Y]=\mu_X+\frac{\mathrm{Cov}(X,Y)}{\mathrm{Var}(Y)}(Y-\mu_Y).
  $$
  \item Show that
  $$
  \mathrm{Var}(X\mid Y)=\sigma_X^2(1-\rho^2).
  $$
  \item (Finance) Interpret $\rho$ as ``hedge effectiveness'' and relate $\mathrm{Var}(X\mid Y)$ to residual risk.
\end{enumerate}
\end{exercise}

% =====================================================
%  Exercise 4
% =====================================================
\begin{exercise}[One-step binomial pricing (mini finance)]
A stock price satisfies $S_0=100$. After one period,
$$
S_1=
\begin{cases}
120 & \text{(up state)},\\
90  & \text{(down state)}.
\end{cases}
$$
The \emph{physical} probability of the up state is $p\in(0,1)$.
We price by \emph{no-arbitrage}. The risk-free rate is $r=0$,
so the bank account satisfies $B_0=1$ and $B_1=1$.

Consider a European call option with strike $K=100$:
$$
H=(S_1-K)_+.
$$

\begin{enumerate}[label=\textbf{(\alph*)}, itemsep=0.35em]

  \item \textbf{Payoff table.}
  Compute the option payoff $H$ in each state and fill:
  $$
  \begin{array}{c|c|c}
  \text{State} & S_1 & H=(S_1-100)_+\\ \hline
  \text{Up} & 120 & \ ? \ \\
  \text{Down} & 90 & \ ? \ \\
  \end{array}
  $$

  \item \textbf{Risk-neutral probability.}
  A probability $q\in(0,1)$ is \emph{risk-neutral} if the discounted stock price
  is a martingale. Since $r=0$, this means
  $$
  \mathbb E^{q}[S_1]=S_0.
  $$
  \begin{enumerate}[label=\textbf{(\roman*)}, itemsep=0.2em]
    \item Write $\mathbb E^{q}[S_1]$ explicitly in terms of $q$.
    \item Solve for $q$ and check that $0<q<1$.
    \item Explain in one sentence why $q$ is not necessarily equal to the real probability $p$.
  \end{enumerate}

  \item \textbf{No-arbitrage option price.}
  The no-arbitrage price is the risk-neutral expectation (no discounting here):
  $$
  V_0=\mathbb E^{q}[H].
  $$
  \begin{enumerate}[label=\textbf{(\roman*)}, itemsep=0.2em]
    \item Write $\mathbb E^{q}[H]$ using the two-state payoff table.
    \item Compute $V_0$ numerically.
  \end{enumerate}

  \item \textbf{Replicating portfolio.}
  Find $\Delta$ (shares) and $B$ (cash) such that
  $$
  \Delta S_1 + B = H
  \quad\text{in both states.}
  $$
  \begin{enumerate}[label=\textbf{(\roman*)}, itemsep=0.2em]
    \item Write the two equations and solve for $\Delta$ and $B$.
    \item Compute the initial cost $\Delta S_0 + B$ and check it matches $V_0$.
  \end{enumerate}

  \item \textbf{Short interpretation.}
  \begin{enumerate}[label=\textbf{(\roman*)}, itemsep=0.2em]
    \item Is $B$ positive or negative? Interpret (lending vs borrowing).
    \item Is $\Delta$ between $0$ and $1$? Interpret as a hedge ratio.
  \end{enumerate}

\end{enumerate}
\end{exercise}



% ==========================================================
\section{Lecture 2 --- Option Pricing : Brownian Motion and Quadratic Variation.}

\begin{exercise}[Gaussian increments \& scaling warm-up]
Let $0\le s<t$.
\begin{enumerate}
\item Compute the law of $B_t-B_s$ and its mean/variance.
\item Show that $\mathrm{Cov}(B_s,B_t)=s$ and $\E[B_sB_t]=s$.
\item (Scaling) For $c>0$, define $W_t:=\frac{1}{\sqrt c}B_{ct}$.
Show that $(W_t)_{t\ge 0}$ is again a standard Brownian motion.
\end{enumerate}
\end{exercise}

% ----------------------------------------------------------
\begin{exercise}[$B_t^2-t$ is a martingale]
\begin{enumerate}
\item Using It\^o's formula, compute $\mathrm{d}(B_t^2)$.
\item Deduce that $B_t^2-t$ is a martingale.
\item Use this to compute $\E[B_t^2]$.
\end{enumerate}
\end{exercise}

% ----------------------------------------------------------
\begin{exercise}[Exponential martingale (drift killing)]
Fix $b\in\R$ and define
$$
Y_t:=\exp\Big(bB_t-\tfrac12 b^2 t\Big).
$$
\begin{enumerate}
\item Compute $\mathrm{d} Y_t$ using It\^o and show that $(Y_t)$ is a martingale.
\item Deduce $\E[e^{bB_t}]=e^{\frac12 b^2 t}$.
\item (Optional) Show that $\E[Y_t^p]<\infty$ for every $p>0$ and compute it.
\end{enumerate}
\end{exercise}

% ----------------------------------------------------------
\begin{exercise}[Hitting a level: reflection principle]
Fix $a>0$ and define the first passage time $\tau_a:=\inf\{t\ge 0:\,B_t=a\}$.
\begin{enumerate}
\item Show that for every $t>0$,
$$
\mathbb{P}(\tau_a\le t)=2\,\mathbb{P}(B_t\ge a).
$$
\item Deduce an explicit formula for $\mathbb{P}(\tau_a\le t)$ in terms of the Gaussian cdf.
\item Differentiate (formally) to guess the density of $\tau_a$.
\end{enumerate}
\end{exercise}

% ----------------------------------------------------------
\begin{exercise}[Max distribution]
Let $M_t:=\sup_{0\le s\le t}B_s$.
\begin{enumerate}
\item Show that for $m>0$, $\mathbb{P}(M_t\ge m)=2\,\mathbb{P}(B_t\ge m)$.
\item Deduce the cdf of $M_t$ and compute $\E[M_t]$ (hint: integrate the tail).
\end{enumerate}
\end{exercise}

% ----------------------------------------------------------
\begin{exercise}[Joint law of $(M_t,B_t)$ (application)]
Assume you are given that $(M_t,B_t)$ has density, for $m>0$ and $w\le m$,
$$
f_{M_t,B_t}(m,w)=\frac{2(2m-w)}{t\sqrt{2\pi t}}
\exp\!\Big(-\frac{(2m-w)^2}{2t}\Big).
$$
\begin{enumerate}
\item Check that $f_{M_t,B_t}$ integrates to $1$.
\item Compute the conditional density of $M_t$ given $B_t=w$ (for $w\in\R$).
\item Compute $\mathbb{P}(M_t\le m\,|\,B_t=w)$ for $m\ge \max(w,0)$.
\end{enumerate}
\end{exercise}

% ----------------------------------------------------------
\begin{exercise}[Barrier probability from reflection]
Let $H>0$. Compute explicitly
$$
\mathbb{P}(M_t < H)\qquad\text{and}\qquad \mathbb{P}(M_t \ge H).
$$
Then interpret $\{M_t<H\}$ as a “survival” event for an up-and-out barrier.
\end{exercise}

% ----------------------------------------------------------
% ----------------------------------------------------------
\begin{exercise}[Exit time of an interval $(a,b)$ with $a<0<b$ (guided)]
Let $(B_t)_{t\ge 0}$ be a standard Brownian motion with $B_0=0$ and fix
$$
a<0<b.
$$
Define the exit time
$$
\tau:=\inf\{t\ge 0:\, B_t\notin(a,b)\}.
$$

\begin{enumerate}[label=\textbf{(\alph*)}, itemsep=0.35em]

\item \textbf{Where can you exit?}
Using the continuity of Brownian paths, show that
$$
B_\tau \in \{a,b\}\quad\text{a.s.}
$$

\item \textbf{Truncation (bounded stopping times).}
For $n\in\mathbb{N}$ define the bounded stopping time
$$
\tau_n:=\tau\wedge n.
$$
Check that $\tau_n\le n$ and that $B_{\tau_n}$ is integrable.

(\emph{Hint:} argue that $B_{t\wedge\tau}\in[a,b]$ for all $t\ge 0$, hence $|B_{\tau_n}|\le \max\{|a|,|b|\}$.)

\item \textbf{Hitting probabilities via the martingale $B_t$.}
Let
$$
p:=\mathbb{P}(B_\tau=b).
$$
\begin{enumerate}[label=\roman*)]
\item Apply optional stopping to the martingale $(B_t)$ at $\tau_n$:
$$
\E[B_{\tau_n}]=0.
$$
\item Justify briefly that you may pass to the limit $n\to\infty$ and obtain
$$
\E[B_\tau]=0.
$$
(\emph{Hint:} dominated convergence using the bound $|B_{\tau_n}|\le \max\{|a|,|b|\}$.)
\item Use $B_\tau\in\{a,b\}$ a.s. to write
$$
\E[B_\tau]=a(1-p)+bp,
$$
and solve for $p$. Deduce $\mathbb{P}(B_\tau=a)=1-p$.
\end{enumerate}

\item \textbf{Expected exit time via the martingale $B_t^2-t$.}
Consider the martingale
$$
M_t:=B_t^2-t.
$$
\begin{enumerate}[label=\roman*)]
\item Apply optional stopping to $(M_t)$ at $\tau_n$:
$$
\E[B_{\tau_n}^2-\tau_n]=0.
$$
\item Rearrange:
$$
\E[\tau_n]=\E[B_{\tau_n}^2].
$$
\item Show that $B_{\tau_n}\to B_\tau$ a.s. and that $|B_{\tau_n}|\le \max\{|a|,|b|\}$, hence
$$
\E[B_{\tau_n}^2]\to \E[B_\tau^2].
$$
(\emph{Hint:} dominated convergence.)
\item Conclude that $\E[\tau_n]\to \E[\tau]$ and obtain
$$
\E[\tau]=\E[B_\tau^2].
$$
Finally compute $\E[B_\tau^2]$ using the probabilities from the previous question and
$$
B_\tau^2=
\begin{cases}
a^2 & \text{if } B_\tau=a,\\
b^2 & \text{if } B_\tau=b.
\end{cases}
$$
\end{enumerate}

\end{enumerate}
\end{exercise}

% ----------------------------------------------------------
% ----------------------------------------------------------
%\begin{exercise}[Laplace transform of exit time (medium)]
%With the same $(\tau,1,b,x)$ as above, fix $\lambda>0$ and set
%$$
%u(x):=\E_x\!\left[e^{-\lambda \tau}\,;\,B_\tau=b\right].
%$$
%\begin{enumerate}
%\item Argue (heuristically or using Markov property) that $u$ solves the ODE
%$$
%\frac12 u''(x)=\lambda u(x),\qquad x\in(1,b),
%$$
%with boundary conditions $u(1)=0$, $u(b)=1$.
%\item Solve the ODE and give $u(x)$ explicitly.
%\end{enumerate}
%\end{exercise}

% ----------------------------------------------------------
\begin{exercise}[A clean identity: $\int_0^t f(s)\,\mathrm{d} B_s$ without stochastic integral]
Let $f:[0,T]\to\R$ be $C^1$.
\begin{enumerate}
\item Apply It\^o to $f(t)B_t$ and show that
$$
\int_0^t f(s)\,\mathrm{d} B_s = f(t)B_t-\int_0^t f'(s)B_s\,\mathrm{d} s.
$$
\item Apply this to $f(s)=e^{\alpha s}$ and simplify.
\end{enumerate}
\end{exercise}

% ----------------------------------------------------------
\begin{exercise}[Compute $\int_0^t e^{B_s}\,\mathrm{d} B_s$ without $\mathrm{d} B$]
\begin{enumerate}
\item Apply It\^o to $e^{B_t}$ and isolate the term $\int_0^t e^{B_s}\,\mathrm{d} B_s$.
\item Deduce an expression for $\int_0^t e^{B_s}\,\mathrm{d} B_s$ involving only $e^{B_t}$ and ordinary integrals.
\end{enumerate}
\end{exercise}

% ----------------------------------------------------------
\begin{exercise}[Quadratic variation on a uniform grid]
Let $0=t_0<\dots<t_n=T$ be the uniform grid $t_k=kT/n$ and set
$$
Q_T^{(n)}:=\sum_{k=0}^{n-1}\big(B_{t_{k+1}}-B_{t_k}\big)^2.
$$
\begin{enumerate}
\item Compute $\E[Q_T^{(n)}]$ and $\mathrm{Var}(Q_T^{(n)})$.
\item Deduce that $Q_T^{(n)}\to T$ in $L^2$.
\item Explain briefly why this suggests that Riemann--Stieltjes integrals are unstable for BM.
\end{enumerate}
\end{exercise}

% ----------------------------------------------------------
\begin{exercise}[Left vs right sums for $\phi_t=B_t$]
On the uniform grid $t_k=kT/n$, define
$$
L_n:=\sum_{k=0}^{n-1} B_{t_k}\,(B_{t_{k+1}}-B_{t_k}),
\qquad
R_n:=\sum_{k=0}^{n-1} B_{t_{k+1}}\,(B_{t_{k+1}}-B_{t_k}).
$$
\begin{enumerate}
\item Show that $R_n-L_n=\sum_{k=0}^{n-1}(B_{t_{k+1}}-B_{t_k})^2=Q_T^{(n)}$.
\item Deduce that $R_n-L_n\to T$ in probability.
\item Use the algebraic identity
$$
B_{t_{k+1}}^2-B_{t_k}^2 = 2B_{t_k}(B_{t_{k+1}}-B_{t_k})+(B_{t_{k+1}}-B_{t_k})^2
$$
to guess the limit of $L_n$.
\end{enumerate}
\end{exercise}

% ==========================================================
\section{Lecture 3 --- Stochastic Integral.}

\begin{exercise}[Warm-up: definition on simple processes, mean zero, isometry]
Let $(W_t)_{t\ge 0}$ be a 1D Brownian motion with its natural filtration $(\F_t)$.
Fix $T>0$ and consider a \emph{simple} (step) process on $[0,T]$:
$$
\phi_t=\sum_{k=0}^{n-1}\phi_{t_k}\,\1_{(t_k,t_{k+1}]}(t),
\qquad 0=t_0<\cdots<t_n=T,\qquad \phi_{t_k}\in\F_{t_k}.
$$
Define the It\^o integral by the left-point sum
$$
\int_0^T \phi_t\,\mathrm{d} W_t
:=\sum_{k=0}^{n-1}\phi_{t_k}\big(W_{t_{k+1}}-W_{t_k}\big).
$$
\begin{enumerate}[label=\textbf{(\alph*)}, itemsep=0.25em]
\item Prove that $\E\!\left[\int_0^T \phi_t\,\mathrm{d} W_t\right]=0$.
\item Prove the It\^o isometry for simple processes:
$$
\E\!\left[\left(\int_0^T \phi_t\,\mathrm{d} W_t\right)^2\right]
=\E\!\left[\int_0^T \phi_t^2\,\mathrm{d} t\right]
=\sum_{k=0}^{n-1}\E[\phi_{t_k}^2]\,(t_{k+1}-t_k).
$$
\item (Bonus) Prove the polarized identity for two simple processes $\phi,\psi$:
$$
\E\!\left[\Big(\int_0^T \phi_t\,\mathrm{d} W_t\Big)\Big(\int_0^T \psi_t\,\mathrm{d} W_t\Big)\right]
=\E\!\left[\int_0^T \phi_t\psi_t\,\mathrm{d} t\right].
$$
\end{enumerate}

\smallskip
\noindent\textit{Hint.}
Use conditional expectation w.r.t.\ $\F_{t_k}$ and the independence/centering of increments.
For (b), expand the square and show cross-terms vanish by conditioning at the later time.
\end{exercise}

% ----------------------------------------------------------
\begin{exercise}[Why the left-point matters: right sums for $\phi_t=W_t$]
Fix $t>0$ and a partition $\pi=\{0=t_0<\cdots<t_n=t\}$.
Define
$$
S^{\text{left}}_\pi:=\sum_{k=0}^{n-1}W_{t_k}\big(W_{t_{k+1}}-W_{t_k}\big),
\qquad
S^{\text{right}}_\pi:=\sum_{k=0}^{n-1}W_{t_{k+1}}\big(W_{t_{k+1}}-W_{t_k}\big).
$$
\begin{enumerate}[label=\textbf{(\alph*)}, itemsep=0.25em]
\item Show the identity
$$
S^{\text{right}}_\pi - S^{\text{left}}_\pi
=\sum_{k=0}^{n-1}\big(W_{t_{k+1}}-W_{t_k}\big)^2.
$$
\item Assume you know (from Lecture 1--2) that the quadratic variation along refining partitions satisfies
$\sum_{k} (W_{t_{k+1}}-W_{t_k})^2 \to t$ in probability.
Deduce that
$$
\lim_{|\pi|\to 0} S^{\text{right}}_\pi
=
\left(\lim_{|\pi|\to 0} S^{\text{left}}_\pi\right)+t,
$$
whenever the left limit exists.
\item Using the algebraic identity
$W_{t_{k+1}}^2 - W_{t_k}^2
= 2W_{t_k}(W_{t_{k+1}}-W_{t_k}) + (W_{t_{k+1}}-W_{t_k})^2$,
show that
$$
\int_0^t W_s\,\mathrm{d} W_s
=\frac12\big(W_t^2 - t\big).
$$
\end{enumerate}

\smallskip
\noindent\textit{Comment.}
This exercise is precisely the “pathology” motivating the left-point convention. % cf slides
\end{exercise}

% ----------------------------------------------------------
\begin{exercise}[Deterministic integrand: Gaussianity + variance]
Let $f\in L^2([0,T])$ be deterministic.
\begin{enumerate}[label=\textbf{(\alph*)}, itemsep=0.25em]
\item First assume that $f$ is a step function:
$f(t)=\sum_{k=0}^{n-1}c_k\,\1_{(t_k,t_{k+1}]}(t)$.
Show that
$$
\int_0^T f(t)\,\mathrm{d} W_t
=\sum_{k=0}^{n-1}c_k\big(W_{t_{k+1}}-W_{t_k}\big)
$$
is Gaussian, centered, with variance $\int_0^T f(t)^2\,\mathrm{d} t$.
\item Deduce (by approximation in $L^2$ and the isometry) that for general deterministic
$f\in L^2([0,T])$,
$$
\int_0^T f(t)\,\mathrm{d} W_t \sim \mathcal{N}\!\left(0,\int_0^T f(t)^2\,\mathrm{d} t\right).
$$
\item Compute explicitly the moment generating function
$\E[\exp(\lambda\int_0^T f(t)\,\mathrm{d} W_t)]$ for $\lambda\in\R$.
\end{enumerate}

\smallskip
\noindent\textit{Hint.}
A finite linear combination of independent Gaussians is Gaussian; for (c) use the mgf of
$\mathcal{N}(0,\sigma^2)$.
\end{exercise}

% ----------------------------------------------------------
\begin{exercise}[A Brownian motion built from another one (orthogonal integrand)]
Let $W$ be a Brownian motion in $\R^d$.
Let $A:[0,\infty)\to\R^{d\times d}$ be deterministic, measurable, and assume
$$
A(t)A(t)^\top = I_d \qquad \text{for all } t\ge 0.
$$
Define the $\R^d$-valued process
$$
B_t := \int_0^t A(s)\,\mathrm{d} W_s.
$$
\begin{enumerate}[label=\textbf{(\alph*)}, itemsep=0.25em]
\item Show that $B$ has centered Gaussian increments.
\item Compute the covariance matrix of $B_t-B_s$ for $0\le s<t$.
\item Deduce that $B$ is a Brownian motion in $\R^d$.
\end{enumerate}

\smallskip
\noindent\textit{Hint.}
First check the claim for step-matrix functions $A$, then extend by $L^2$ approximation
using the vector It\^o isometry.
\end{exercise}


% ==========================================================
\section{Lecture 4 --- Ito Calculus and Formula.}

\begin{exercise}[Warm-up: It\^o formula for $f(t,W_t)$]
Let $(W_t)_{t\ge 0}$ be a standard Brownian motion and let $f\in C^{1,2}([0,T]\times\R)$.
\begin{enumerate}[label=\textbf{(\alph*)}, itemsep=0.25em]
  \item State It\^o's formula for $f(t,W_t)$ in differential form and in integral form.
  \item Apply it to $f(t,x)=x^2$ and deduce the decomposition
  $$
    W_t^2 = t + 2\int_0^t W_s\,\mathrm{d} W_s.
  $$
  \item Apply it to $f(t,x)=x^3$ and write $W_t^3$ as the sum of a Lebesgue integral and an It\^o integral.
\end{enumerate}
\end{exercise}

\begin{exercise}[Moments without Gaussian computations]
Let $(W_t)$ be a Brownian motion and fix $T>0$.
\begin{enumerate}[label=\textbf{(\alph*)}, itemsep=0.25em]
  \item Apply It\^o's formula to $f(x)=x^4$ and write $W_T^4$ as
  $$
    W_T^4 = \int_0^T a_s\,\mathrm{d} s + \int_0^T b_s\,\mathrm{d} W_s
  $$
  for explicit adapted processes $(a_s)$ and $(b_s)$.
  \item Take expectations and use $\E[W_t^2]=t$ to recover $\E[W_T^4]=3T^2$.
  \item (Optional) Repeat with $f(x)=x^6$ and obtain a closed form for $\E[W_T^6]$.
\end{enumerate}
\end{exercise}


\begin{exercise}[Exponential Brownian martingale]
Fix $b\in\R$ and define
$$
Y_t := \exp\!\big(a t + b W_t\big), \qquad t\in[0,T],
$$
where $a\in\R$ is to be chosen.
\begin{enumerate}[label=\textbf{(\alph*)}, itemsep=0.25em]
  \item Compute $\mathrm{d} Y_t$ using It\^o's formula.
  \item Choose $a=a(b)$ so that the \emph{drift term} vanishes.
  \item Conclude that $(Y_t)_{0\le t\le T}$ is a martingale, and deduce $\E[Y_T]$.
\end{enumerate}
\end{exercise}

\begin{exercise}[Generalized geometric Brownian motion]
Let $X$ be the It\^o process
$$
\mathrm{d} X_t = \alpha_t\,\mathrm{d} t + \sigma_t\,\mathrm{d} W_t,\qquad X_0=x_0,
$$
where $(\alpha_t)$ and $(\sigma_t)$ are adapted and satisfy
$$
\E\!\left[\int_0^T |\alpha_t|\,\mathrm{d} t\right]<\infty,
\qquad
\E\!\left[\int_0^T \sigma_t^2\,\mathrm{d} t\right]<\infty.
$$
Define $S_t := \exp(X_t)$.
\begin{enumerate}[label=\textbf{(\alph*)}, itemsep=0.25em]
  \item Apply It\^o's formula to $f(x)=e^x$ and derive the dynamics of $(S_t)$:
  write $\mathrm{d} S_t$ in the form $\mathrm{d} S_t = S_t(\cdots)\mathrm{d} t + S_t(\cdots)\mathrm{d} W_t$.
  \item Specialize to $\alpha_t = r_t - \tfrac12\sigma_t^2$ and show that
  $$
    \mathrm{d} S_t = r_t S_t\,\mathrm{d} t + \sigma_t S_t\,\mathrm{d} W_t.
  $$
  \item (Optional) Assume $r_t$ is deterministic and integrable, and set
  $$
    \widetilde S_t := \exp\!\Big(-\int_0^t r_u\,\mathrm{d} u\Big)S_t.
  $$
  Compute $\mathrm{d}\widetilde S_t$ and give a condition ensuring $(\widetilde S_t)$ is a martingale.
\end{enumerate}
\end{exercise}


\begin{exercise}[Product and ratio of two It\^o processes]
Let $X$ and $Y$ satisfy
$$
\mathrm{d} X_t = \mu^X_t\,\mathrm{d} t + \sigma^X_t\,\mathrm{d} W_t,
\qquad
\mathrm{d} Y_t = \mu^Y_t\,\mathrm{d} t + \sigma^Y_t\,\mathrm{d} W_t,
$$
where all coefficients are adapted and square-integrable on $[0,T]$ (as needed).
\begin{enumerate}[label=\textbf{(\alph*)}, itemsep=0.25em]
  \item Set $Z_t:=X_tY_t$. Using It\^o's product rule, derive the dynamics of $Z_t$
  (identify both the drift and diffusion terms explicitly).
  \item Assume now that $X_t>0$ and $Y_t>0$ a.s.\ for all $t$. Set $R_t:=X_t/Y_t$.
  Compute $\mathrm{d}(\ln R_t)$ using It\^o's formula, then deduce the SDE for $R_t$.
  \item (Optional) Consider the multiplicative dynamics
  $$
    \mathrm{d} X_t = X_t(\mu^X_t\,\mathrm{d} t + \sigma^X_t\,\mathrm{d} W_t),
    \qquad
    \mathrm{d} Y_t = Y_t(\mu^Y_t\,\mathrm{d} t + \sigma^Y_t\,\mathrm{d} W_t).
  $$
  Solve explicitly for $X_t$ and $Y_t$ (as stochastic exponentials), and verify that the resulting expression
  for $R_t=X_t/Y_t$ matches your answer in (b).
\end{enumerate}
\end{exercise}

\section{Lecture 5 --- Black--Scholes price and Greeks.}


Throughout, fix $0\le t<T$, $\tau:=T-t$, $K>0$, $r\in\R$, $\sigma>0$, and denote by
$$
\Phi(x):=\int_{-\infty}^x \frac{1}{\sqrt{2\pi}}e^{-u^2/2}\,du,
\qquad
\varphi(x):=\Phi'(x)=\frac{1}{\sqrt{2\pi}}e^{-x^2/2}
$$
the standard normal CDF and density.

We also use the usual notations
$$
d_1(t,s):=\frac{\ln(s/K)+(r+\frac12\sigma^2)\tau}{\sigma\sqrt{\tau}},
\qquad
d_2(t,s):=d_1(t,s)-\sigma\sqrt{\tau}.
$$
The Black--Scholes call and put prices are
$$
C(t,s)=s\Phi(d_1)-K e^{-r\tau}\Phi(d_2),
\qquad
P(t,s)=K e^{-r\tau}\Phi(-d_2)-s\Phi(-d_1).
$$

% ----------------------------------------------------------
\subsection*{Warm-up identities (very useful)}
% ----------------------------------------------------------

\begin{exercise}[Derivatives of $d_1,d_2$]
Show that for $s>0$,
$$
\partial_s d_1(t,s)=\partial_s d_2(t,s)=\frac{1}{s\sigma\sqrt{\tau}}.
$$
Compute also $\partial_\sigma d_1(t,s)$ and $\partial_\sigma d_2(t,s)$.
\end{exercise}

\begin{exercise}[The key ``$\varphi$-identity'']
Prove the identity
\begin{equation}\label{eq:phi-identity-exo}
s\,\varphi(d_1)=K e^{-r\tau}\,\varphi(d_2).
\end{equation}
\emph{Hint:} write $\varphi(d_2)/\varphi(d_1)=\exp(-(d_2^2-d_1^2)/2)$ and use
$d_2=d_1-\sigma\sqrt{\tau}$.
\end{exercise}

\begin{exercise}[Put--call parity]
Prove
$$
C(t,s)-P(t,s)=s-K e^{-r\tau}.
$$
Then deduce $\Delta_P=\Delta_C-1$ and $\Gamma_P=\Gamma_C$ once you have computed Greeks for the call.
\end{exercise}

% ----------------------------------------------------------
\subsection*{Greeks: compute and prove}
% ----------------------------------------------------------

\begin{exercise}[Delta of the call]
Starting from
$$
C(t,s)=s\Phi(d_1)-K e^{-r\tau}\Phi(d_2),
$$
differentiate w.r.t.\ $s$ and prove
$$
\boxed{\ \Delta(t,s):=\partial_s C(t,s)=\Phi(d_1(t,s)).\ }
$$
\emph{Hint:} you will obtain extra terms involving $\varphi(d_1)\partial_s d_1$ and
$\varphi(d_2)\partial_s d_2$; use \eqref{eq:phi-identity-exo} to cancel them.
\end{exercise}

\begin{exercise}[Gamma of the call]
Differentiate your Delta formula and show
$$
\boxed{\ \Gamma(t,s):=\partial_{ss}C(t,s)=\frac{\varphi(d_1(t,s))}{s\sigma\sqrt{\tau}}.\ }
$$
Then show $\Gamma(t,s)>0$ for all $s>0$.
\end{exercise}

\begin{exercise}[Vega of the call]
Differentiate $C(t,s)$ with respect to $\sigma$ and prove
$$
\boxed{\ \text{Vega}(t,s):=\partial_\sigma C(t,s)= s\,\varphi(d_1(t,s))\,\sqrt{\tau}.\ }
$$
Conclude that the call price is strictly increasing in $\sigma$ (for $\tau>0$).
\end{exercise}

\begin{exercise}[Rho of the call]
Differentiate with respect to $r$ and show
$$
\boxed{\ \rho(t,s):=\partial_r C(t,s)=K \tau e^{-r\tau}\Phi(d_2(t,s)).\ }
$$
\emph{Optional:} interpret the sign of $\rho$ for calls and puts.
\end{exercise}

\begin{exercise}[Theta of the call]
Compute $\Theta(t,s):=\partial_t C(t,s)$ and prove the formula
$$
\boxed{\ 
\Theta(t,s)= -\frac{s\varphi(d_1)\sigma}{2\sqrt{\tau}}-rK e^{-r\tau}\Phi(d_2).
\ }
$$
\emph{Hint:} it is often simpler to differentiate in $\tau=T-t$ first, then use
$\partial_t=-\partial_\tau$.
\end{exercise}

\begin{exercise}[A Greek identity from the PDE]
Assume (or recall) that $C$ satisfies the Black--Scholes PDE
$$
\partial_t C + rs\,\partial_s C + \tfrac12\sigma^2 s^2 \partial_{ss}C - rC = 0.
$$
Rewrite it as a relation between Greeks:
$$
\boxed{\ \Theta + rs\,\Delta + \tfrac12\sigma^2 s^2 \Gamma - rC = 0.\ }
$$
Check the identity by plugging your closed-form expressions.
\end{exercise}

% ----------------------------------------------------------
\subsection*{Conceptual / interpretation exercises}
% ----------------------------------------------------------

\begin{exercise}[Qualitative behavior of Delta]
Show that for a call, $\Delta(t,s)\in(0,1)$ and that
$$
\Delta(t,s)\to 1\ \text{as } s\to\infty,
\qquad
\Delta(t,s)\to 0\ \text{as } s\downarrow 0.
$$
Interpret these limits in terms of hedging.
\end{exercise}

\begin{exercise}[Convexity and Gamma]
Let $g(s)=(s-K)_+$. Show that $g$ is convex and explain why, in a diffusion model,
the option price $s\mapsto C(t,s)$ should be convex (hence $\Gamma\ge 0$).
\emph{Hint:} you may use Jensen's inequality on the risk-neutral pricing formula if you have seen it,
or argue directly from the explicit $\Gamma$ formula.
\end{exercise}

\begin{exercise}[Short-maturity behavior (ATM scaling)]
Assume $s=K$ and $r=0$ for simplicity. Study the asymptotic behavior of
$$
C(t,K) \quad \text{as }\ \tau\downarrow 0.
$$
Show that $C(t,K)\sim K\,\sigma\sqrt{\tau}\,\varphi(0)$.
Interpret this as ``time value is of order $\sqrt{\tau}$''.
\end{exercise}

% ----------------------------------------------------------
\subsection*{Implied volatility (optional but very standard)}
% ----------------------------------------------------------

\begin{exercise}[No-arbitrage bounds]
Show that for a call,
$$
(s-K e^{-r\tau})_+ \le C(t,s)\le s.
$$
\emph{Hint:} use $0\le (S_T-K)_+\le S_T$ and take discounted expectations.
\end{exercise}

\begin{exercise}[Existence/uniqueness of implied vol (using Vega)]
Fix $(t,s,K,T,r)$ and suppose you observe a price $C^{\mathrm{mkt}}$ satisfying the
bounds of the previous exercise and $\tau>0$.
Assuming continuity of $\sigma\mapsto C(t,s;\sigma)$, use Vega $>0$ to argue that the
implied volatility equation
$$
C(t,s;K,T,r,\sigma_{\mathrm{imp}})=C^{\mathrm{mkt}}
$$
has at most one solution. (Discuss briefly why a solution should exist.)
\end{exercise}

\begin{exercise}[Implied-vol sensitivity (Vega inversion)]
Assume $\sigma_{\mathrm{imp}}$ is defined implicitly by
$C(t,s;\sigma_{\mathrm{imp}})=C^{\mathrm{mkt}}$.
Differentiate both sides w.r.t.\ $C^{\mathrm{mkt}}$ and show
$$
\boxed{\ \frac{\partial \sigma_{\mathrm{imp}}}{\partial C^{\mathrm{mkt}}}
=
\frac{1}{\text{Vega}(t,s;K,T,r,\sigma_{\mathrm{imp}})}.\ }
$$
\emph{Interpretation:} when Vega is small (deep ITM/OTM, short maturity), implied vol is unstable.
\end{exercise}


% % --- Put ~10--15 exercises here ---

\section{Lecture 6 \& 7 --- Girsanov Theorem and Neutral Risk Measure.}

% ============================================================
% Extra exercises (medium -> hard) aligned with the notes



% ------------------------------------------------------------
\begin{exercise}[Bayes rule in practice: conditional expectation under $\mathbb{Q}$.]
Let $Z\ge 0$ with $\E^\mathbb{P}[Z]=1$ and define $\mathbb{Q}$ by $\mathrm{d}\mathbb{Q}=Z\,\mathrm{d}\mathbb{P}$ on $\F_T$.
Let $0\le s\le t\le T$ and let $Y$ be integrable and $\F_t$-measurable.
\begin{enumerate}[leftmargin=1.2em, label=(\alph*), itemsep=0.3em]
\item Prove that for any $G\in\F_s$,
$$
\E^\mathbb{Q}[Y\1_G]=\E^\mathbb{P}[Z_t\,Y\,\1_G],
$$
where $(Z_u)_{u\le T}$ is the density process $Z_u:=\E^\mathbb{P}[Z_T\mid \F_u]$.
\item Deduce the Bayes formula
$$
\E^\mathbb{Q}[Y\mid \F_s]=\frac{1}{Z_s}\,\E^\mathbb{P}[Z_tY\mid \F_s].
$$
\item Apply this with $Y=\1_{\{W_t\ge a\}}$ and $Z_T=\exp(\theta W_T-\tfrac12\theta^2T)$
(constant tilt) to express $\mathbb{Q}(W_t\ge a\mid \F_s)$ as an explicit function of $(t,s,W_s)$.
\end{enumerate}
\end{exercise}


% ------------------------------------------------------------
\begin{exercise}[Gaussian tilting in dimension $2$ (correlation + conditional law).]
Under $\mathbb{P}$, let $(X,Y)$ be centered Gaussian with
$$
\text{Var}(X)=\text{Var}(Y)=1,\qquad \text{Cov}(X,Y)=\rho\in(-1,1).
$$
Define $\mathbb{Q}$ by $\mathrm{d}\mathbb{Q} = Z\,\mathrm{d}\mathbb{P}$ with
$$
Z=\exp\!\Big(\theta X-\tfrac12\theta^2\Big).
$$
\begin{enumerate}[leftmargin=1.2em, label=(\alph*), itemsep=0.3em]
\item Identify the law of $X$ under $\mathbb{Q}$.
\item Compute $\E^\mathbb{Q}[Y]$ and $\text{Var}^\mathbb{Q}(Y)$.
\item Show that the conditional law of $Y$ given $X$ is the same under $\mathbb{P}$ and $\mathbb{Q}$,
and deduce the joint law of $(X,Y)$ under $\mathbb{Q}$ (mean vector and covariance matrix).
\end{enumerate}
\end{exercise}
% ------------------------------------------------------------
\begin{exercise}[Cameron--Martin check (deterministic drift): characteristic function.]
Let $W$ be a 1D Brownian motion under $\mathbb{P}$ on $[0,T]$ and let $h\in L^2([0,T])$ deterministic.
Define
$$
Z_T=\exp\!\left(\int_0^T h(t)\,\mathrm{d} W_t-\frac12\int_0^T h(t)^2\,\mathrm{d} t\right),
\qquad \frac{\mathrm{d}\mathbb{Q}}{\mathrm{d}\mathbb{P}}\Big|_{\F_T}=Z_T.
$$
Let $B_t:=W_t-\int_0^t h(u)\,\mathrm{d} u$.
\begin{enumerate}[leftmargin=1.2em, label=(\alph*), itemsep=0.3em]
\item For fixed $0\le s<t\le T$ and $u\in\R$, compute
$\E^\mathbb{Q}[e^{iu(B_t-B_s)}\mid \F_s]$.
\item Deduce that $B_t-B_s\sim \mathcal N(0,t-s)$ under $\mathbb{Q}$ and is independent of $\F_s$.
\item Conclude that $B$ is a Brownian motion under $\mathbb{Q}$.
\end{enumerate}
\end{exercise}
% ------------------------------------------------------------
\begin{exercise}[Novikov condition: easy cases and a counterexample search.]
Let $\phi$ be progressively measurable with $\int_0^T \phi_t^2\,\mathrm{d} t<\infty$ a.s.,
and define the exponential martingale
$$
Z_t=\exp\!\left(-\int_0^t \phi_s\,\mathrm{d} W_s-\frac12\int_0^t \phi_s^2\,\mathrm{d} s\right).
$$
\begin{enumerate}[leftmargin=1.2em, label=(\alph*), itemsep=0.3em]
\item Show that if $\phi$ is bounded, then Novikov holds and $(Z_t)$ is a true martingale.
\item Show that if $\phi_t=\alpha W_t$ for a constant $\alpha\in\R$, then
Novikov holds for $T$ small enough (give an explicit condition on $(\alpha,T)$).
\item (Hard) Look for a progressively measurable $\phi$ such that $\int_0^T \phi_t^2\,\mathrm{d} t<\infty$ a.s.
but $\E[Z_T]<1$ (strict supermartingale). Give either a construction or a literature reference.
\end{enumerate}
\end{exercise}
% ------------------------------------------------------------
\begin{exercise}[Girsanov + SDE drift switch: verify on an explicit process.]
Let $W$ be Brownian under $\mathbb{P}$ and let $\phi$ satisfy Novikov.
Define $\mathbb{Q}$ by $\mathrm{d}\mathbb{Q}/\mathrm{d}\mathbb{P}|_{\F_T}=Z_T$ with
$$
Z_T=\exp\!\left(-\int_0^T \phi_s\,\mathrm{d} W_s-\frac12\int_0^T \phi_s^2\,\mathrm{d} s\right),
$$
and let $B_t:=W_t+\int_0^t \phi_s\,\mathrm{d} s$ (Brownian under $\mathbb{Q}$).
Consider the It\^o process under $\mathbb{P}$
$$
\mathrm{d} X_t = \mu_t\,\mathrm{d} t + \sigma_t\,\mathrm{d} W_t,\qquad X_0=x_0,
$$
with adapted coefficients and $\sigma_t\neq 0$.
\begin{enumerate}[leftmargin=1.2em, label=(\alph*), itemsep=0.3em]
\item Rewrite the dynamics of $X$ under $\mathbb{Q}$ in terms of $B$.
\item Choose $\phi_t=-\mu_t/\sigma_t$ and deduce that $X$ is a (local) martingale under $\mathbb{Q}$.
\item Apply to the discounted stock $\bar S_t=D_tS_t$ from Part D with
$\phi_t=(\mu_t-r_t)/\sigma_t$ and recover $\mathrm{d}\bar S_t=\sigma_t \bar S_t\,\mathrm{d} B_t$ under $\mathbb{Q}$.
\end{enumerate}
\end{exercise}
% ------------------------------------------------------------
\begin{exercise}[Risk-neutral pricing: verify the martingale property of discounted option prices.]
Assume constant $(r,\sigma)$ and under $\mathbb{Q}$:
$$
\frac{\mathrm{d} S_t}{S_t}=r\,\mathrm{d} t+\sigma\,\mathrm{d} B_t.
$$
Let $H=g(S_T)$ with $g$ such that $g(S_T)\in L^1(\mathbb{Q})$ and define
$$
V_t:=\E^\mathbb{Q}\!\left[e^{-r(T-t)}H\mid \F_t\right].
$$
\begin{enumerate}[leftmargin=1.2em, label=(\alph*), itemsep=0.3em]
\item Show that the discounted process $\bar V_t:=e^{-rt}V_t$ is a $\mathbb{Q}$-martingale.
\item Use It\^o (assume $V_t=p(t,S_t)$ smooth) to show that $p$ solves
$$
\partial_t p+rs\,\partial_s p+\tfrac12\sigma^2 s^2\,\partial_{ss}p-rp=0,\qquad p(T,s)=g(s).
$$
\item Conclude that ``PDE pricing'' and ``risk-neutral expectation pricing'' coincide.
\end{enumerate}
\end{exercise}
% ------------------------------------------------------------
\begin{exercise}[Discrete-time delta hedging error: first-order expansion (hard).]
Consider the Black--Scholes model with constant $(r,\sigma)$ and a European payoff $g(S_T)$.
Let $p(t,s)$ be the BS price function (smooth enough) and define the discrete delta hedge on a grid
$0=t_0<\cdots<t_n=T$ by $\Delta_{t_{i-1}}=\partial_s p(t_{i-1},S_{t_{i-1}})$ and discounted wealth
$$
e^{-rT}X_T^{\,n}=p(0,S_0)+\sum_{i=1}^n \Delta_{t_{i-1}}\big(\bar S_{t_i}-\bar S_{t_{i-1}}\big),
\qquad \bar S_t=e^{-rt}S_t.
$$
Let $\mathrm{PL}_T^{\,n}:=X_T^{\,n}-g(S_T)$.
\begin{enumerate}[leftmargin=1.2em, label=(\alph*), itemsep=0.3em]
\item Show (formally) that the continuous-time hedge gives $\mathrm{PL}_T^{\,\infty}=0$.
\item Prove that the leading-order term of the (discounted) hedging error satisfies
$$
e^{-rT}\mathrm{PL}_T^{\,n}\approx \frac12\sum_{i=1}^n
\Gamma_{t_{i-1}}\,(\Delta S_i)^2
\quad\text{with}\quad
\Gamma_t=\partial_{ss}p(t,S_t),\ \Delta S_i=S_{t_i}-S_{t_{i-1}},
$$
up to terms of smaller order in the mesh size.
\item Deduce a heuristic scaling for $\text{Var}(\mathrm{PL}_T^{\,n})$ as the mesh size goes to $0$.
\end{enumerate}
\end{exercise}
% ------------------------------------------------------------
\begin{exercise}[Greeks by differentiation under the expectation (bridge to Girsanov).]
In the Black--Scholes model under $\mathbb{Q}$ with constant $(r,\sigma)$,
let
$$
p(t,s)=\E^\mathbb{Q}\!\left[e^{-r(T-t)}g(S_T)\mid S_t=s\right],
\quad
S_T=s\exp\!\left((r-\tfrac12\sigma^2)(T-t)+\sigma\sqrt{T-t}\,Z\right).
$$
Assume $g$ is Lipschitz with polynomial growth.
\begin{enumerate}[leftmargin=1.2em, label=(\alph*), itemsep=0.3em]
\item Justify (dominated convergence) that $\partial_s p$ can be computed by differentiating inside the expectation,
and express $\Delta=\partial_s p$ as an expectation involving $g$.
\item Do the same for Vega $\partial_\sigma p$ (careful: $\sigma$ appears in mean and variance).
\item Apply to $g(s)=(s-K)_+$ and recover $\Delta_t=\mathbf N(d_+(t,S_t))$.
\end{enumerate}

\end{exercise}


% ------------------------------------------------------------
\begin{exercise}[Asian option (arithmetic average): martingale pricing, state augmentation, and PDE.]
Assume constant $(r,\sigma)$ and under $\mathbb{Q}$:
$$
\frac{\mathrm{d}S_t}{S_t}=r\,\mathrm{d}t+\sigma\,\mathrm{d}B_t,
\qquad S_0>0.
$$
Consider the arithmetic-average Asian call with payoff
$$
H=\left(\frac1T\int_0^T S_u\,\mathrm{d}u-K\right)^+.
$$
Define the running integral
$$
I_t:=\int_0^t S_u\,\mathrm{d}u.
$$

\begin{enumerate}[leftmargin=1.2em, label=(\alph*), itemsep=0.3em]
\item Write the risk-neutral pricing formula for the option price
$$
V_t=\E^\mathbb{Q}\!\left[e^{-r(T-t)}\left(\frac1T\int_0^T S_u\,\mathrm{d}u-K\right)^+\Bigm|\F_t\right].
$$

\item Show that
$$
\int_0^T S_u\,\mathrm{d}u=I_t+\int_t^T S_u\,\mathrm{d}u,
$$
and deduce that the price can be written as a deterministic function of \emph{two} state variables:
$$
V_t=v(t,S_t,I_t).
$$

\item Derive the dynamics of $(S_t,I_t)$ under $\mathbb{Q}$:
$$
\mathrm{d}S_t=rS_t\,\mathrm{d}t+\sigma S_t\,\mathrm{d}B_t,
\qquad
\mathrm{d}I_t=S_t\,\mathrm{d}t.
$$
Explain why $(S_t,I_t)$ is Markov even though $S$ alone is not enough for pricing this payoff.

\item Assume $v\in C^{1,2,1}$ is smooth enough. Apply It\^o's formula to $v(t,S_t,I_t)$ and derive the PDE satisfied by $v$:
$$
\partial_t v + rs\,\partial_s v + s\,\partial_i v + \frac12\sigma^2 s^2\,\partial_{ss}v - rv=0,
$$
with terminal condition
$$
v(T,s,i)=\left(\frac{i}{T}-K\right)^+.
$$

\item Define the discounted price $\bar V_t=e^{-rt}V_t$. Show directly from the risk-neutral pricing formula that $(\bar V_t)$ is a $\mathbb{Q}$-martingale.

\item (Numerical / implementation-oriented) Propose a Monte Carlo estimator for $V_0$ based on an Euler discretization of
$$
\int_0^T S_u\,\mathrm{d}u
\approx \sum_{k=0}^{n-1} S_{t_k}\Delta t.
$$
State clearly:
\begin{itemize}[leftmargin=1.2em]
\item what is simulated under $\mathbb{Q}$,
\item what is averaged,
\item where the discount factor appears.
\end{itemize}

\item (Bonus, medium-hard) Let
$$
G:=\exp\!\left(\frac1T\int_0^T \log S_u\,\mathrm{d}u\right)
$$
(the geometric average). Explain why the geometric-average Asian option is analytically tractable in Black--Scholes, and suggest using it as a \emph{control variate} for Monte Carlo pricing of the arithmetic Asian option.
\end{enumerate}
\end{exercise}
% ------------------------------------------------------------
\begin{exercise}[Change of num\'eraire (forward measure) in a one-step pricing identity.]
Let $B(t,T)$ denote the time-$t$ price of a zero-coupon bond maturing at $T$, and let $N_t:=B(t,T)$ be the num\'eraire.
Assume a measure $\mathbb{Q}^T$ exists such that for any traded asset $X$,
$$
\frac{X_t}{B(t,T)}
$$
is a $\mathbb{Q}^T$-martingale (under suitable integrability assumptions).

Consider a payoff at time $T$ of the form
$$
H=(S_T-K)^+.
$$

\begin{enumerate}[leftmargin=1.2em, label=(\alph*), itemsep=0.3em]
\item Starting from the general num\'eraire pricing formula, show that
$$
V_t = B(t,T)\,\E^{\mathbb{Q}^T}\!\left[\frac{H}{B(T,T)}\Bigm|\F_t\right].
$$

\item Use $B(T,T)=1$ to deduce
$$
V_t = B(t,T)\,\E^{\mathbb{Q}^T}\!\left[H\mid\F_t\right].
$$

\item Compare this with money-market pricing:
$$
V_t=\E^{\mathbb{Q}}\!\left[e^{-\int_t^T r_u\,\mathrm{d}u}H\mid\F_t\right].
$$
Explain in words why both formulas price the \emph{same} claim but under different numeraires/measures.

\item (Conceptual) Give one example from fixed-income where choosing a forward measure is more convenient than using the money-market num\'eraire (e.g.\ caplets, bond options, swaptions).
\end{enumerate}
\end{exercise}
% ------------------------------------------------------------
\begin{exercise}[Incomplete market toy model: many EMMs and non-unique prices.]
Work on a filtered probability space supporting a 2D Brownian motion $(W^1,W^2)$ under $\mathbb{P}$.
Suppose there is only \emph{one} risky traded asset $S$ and one money-market account $B^{(0)}$, with dynamics
$$
\frac{\mathrm{d}S_t}{S_t}=\mu\,\mathrm{d}t+\sigma_1\,\mathrm{d}W_t^1+\sigma_2\,\mathrm{d}W_t^2,
\qquad
\mathrm{d}B_t^{(0)}=rB_t^{(0)}\,\mathrm{d}t,
$$
where $\sigma_1^2+\sigma_2^2>0$.

\begin{enumerate}[leftmargin=1.2em, label=(\alph*), itemsep=0.3em]
\item Write the discounted stock dynamics under $\mathbb{P}$:
$$
\mathrm{d}\bar S_t = \bar S_t\big((\mu-r)\,\mathrm{d}t+\sigma_1\,\mathrm{d}W_t^1+\sigma_2\,\mathrm{d}W_t^2\big).
$$

\item Consider a Girsanov shift $\phi_t=(\phi_t^1,\phi_t^2)$ (for simplicity, take constants $\phi^1,\phi^2$) and define
a measure $\mathbb{Q}$ via the 2D exponential martingale. Show that under $\mathbb{Q}$,
$$
W_t^i = B_t^i - \phi^i t,\qquad i=1,2,
$$
where $(B^1,B^2)$ is a 2D Brownian motion under $\mathbb{Q}$.

\item Derive the condition on $(\phi^1,\phi^2)$ for $\bar S$ to be a local martingale under $\mathbb{Q}$:
$$
\sigma_1\phi^1+\sigma_2\phi^2=\mu-r.
$$

\item Show that this linear equation has infinitely many solutions $(\phi^1,\phi^2)$.
Conclude that the EMM is not unique in this model (market is incomplete).

\item Consider a contingent claim
$$
H=\mathbf{1}_{\{W_T^2\ge 0\}}.
$$
Explain informally why this claim is not replicable using only $(S,B^{(0)})$.

\item (Conceptual pricing consequence) Explain why no-arbitrage does not determine a unique price for $H$:
different EMMs generally give different values of
$$
\E^\mathbb{Q}[e^{-rT}H].
$$
\end{enumerate}
\end{exercise}

% ============================================================
% Exercise --- Continuous-time lookback feature (one-touch)
% ============================================================

\begin{exercise}[Continuous-time lookback digital (one-touch) in Black--Scholes]
We work under the risk-neutral measure $\mathbb{Q}$ (money-market num\'eraire).
Assume the stock pays a continuous dividend yield $q\ge 0$ and follows
\[
\frac{\dd S_t}{S_t}=(r-q)\,\dd t+\sigma\,\dd B_t,
\qquad S_0>0,\ \sigma>0.
\]
Fix a barrier level $H>S_0$ and a maturity $T>0$.

Consider the \emph{lookback digital (one-touch)} payoff
\[
\Pi := \mathbf 1_{\left\{\sup_{0\le t\le T} S_t \ge H\right\}}.
\]
This contract pays \$1 at time $T$ if the stock \emph{ever} touches (or exceeds) $H$ before $T$.

\medskip
\noindent\textbf{Goal.} Compute the no-arbitrage price
\[
V_0 = \E^{\mathbb{Q}}\!\left[e^{-rT}\,\Pi\right]
= e^{-rT}\,\mathbb{Q}\!\left(\sup_{0\le t\le T} S_t \ge H\right).
\]

\begin{enumerate}[leftmargin=1.6em, itemsep=0.35em]

% ------------------------------------------------------------
\item \textbf{Log-transform and reduction to a hitting event.}
Define the log-return process
\[
X_t := \ln\!\Big(\frac{S_t}{S_0}\Big).
\]
\begin{enumerate}[leftmargin=1.6em, itemsep=0.25em]
\item Show (by It\^o's formula) that
\[
X_t = \Big((r-q)-\tfrac12\sigma^2\Big)t + \sigma B_t
=: \mu t + \sigma B_t,
\qquad \mu:=(r-q)-\tfrac12\sigma^2.
\]
\item Show that the barrier event can be rewritten as
\[
\left\{\sup_{0\le t\le T} S_t \ge H\right\}
=
\left\{\sup_{0\le t\le T} X_t \ge a\right\},
\qquad
a:=\ln\!\Big(\frac{H}{S_0}\Big)>0.
\]
\end{enumerate}

% ------------------------------------------------------------
\item \textbf{A key lemma: hitting probability for a drifted Brownian motion.}
Let $(B_t)_{t\ge0}$ be a standard Brownian motion and define
\[
Y_t := \mu t + \sigma B_t,
\qquad \mu\in\R,\ \sigma>0.
\]
Prove the identity
\begin{equation}\label{eq:hitting-drifted}
\mathbb{Q}\!\left(\sup_{0\le t\le T} Y_t \ge a\right)
=
\Phi\!\left(\frac{\mu T-a}{\sigma\sqrt{T}}\right)
+
\exp\!\left(\frac{2\mu a}{\sigma^2}\right)\,
\Phi\!\left(\frac{-\mu T-a}{\sigma\sqrt{T}}\right),
\qquad a>0,
\end{equation}
where $\Phi$ is the standard normal CDF.

\medskip
\noindent\textbf{Guided proof (recommended).}
\begin{enumerate}[leftmargin=1.6em, itemsep=0.25em]
\item First treat the driftless case $\mu=0$:
show that
\[
\mathbb{Q}\!\left(\sup_{0\le t\le T} (\sigma B_t)\ge a\right)
=
2\,\mathbb{Q}(\sigma B_T \ge a)
=
2\left[1-\Phi\!\left(\frac{a}{\sigma\sqrt{T}}\right)\right].
\]
Hint: use the reflection principle for Brownian motion.

\item Now handle $\mu\neq 0$ using a \emph{change of measure}.
Define $\theta:=\mu/\sigma$ and the density
\[
Z_T:=\exp\!\left(-\theta B_T - \tfrac12\theta^2 T\right).
\]
Under the tilted measure $\widetilde{\mathbb{Q}}$ defined by
$\dd\widetilde{\mathbb{Q}}=Z_T\,\dd\mathbb{Q}$ on $\F_T$,
the process $\widetilde B_t:=B_t+\theta t$ is a Brownian motion.
Rewrite $Y_t$ under $\widetilde{\mathbb{Q}}$ in terms of $\widetilde B$ and reduce the problem to
a driftless maximum event under $\widetilde{\mathbb{Q}}$.

\item Carefully transform back to $\mathbb{Q}$ (use $\E^{\mathbb{Q}}[\cdot]=\E^{\widetilde{\mathbb{Q}}}[Z_T^{-1}\cdot]$)
and simplify to obtain \eqref{eq:hitting-drifted}.
\end{enumerate}

% ------------------------------------------------------------
\item \textbf{Closed-form price of the one-touch (lookback digital).}
Use \eqref{eq:hitting-drifted} with $Y=X$, i.e.\ $\mu=(r-q)-\tfrac12\sigma^2$ and $a=\ln(H/S_0)$, to deduce
\[
V_0
=
e^{-rT}\left[
\Phi\!\left(\frac{\mu T-a}{\sigma\sqrt{T}}\right)
+
\exp\!\left(\frac{2\mu a}{\sigma^2}\right)\,
\Phi\!\left(\frac{-\mu T-a}{\sigma\sqrt{T}}\right)
\right],
\qquad
\mu=(r-q)-\tfrac12\sigma^2,\ a=\ln\!\Big(\frac{H}{S_0}\Big).
\]
\textbf{Sanity checks.}
\begin{itemize}[leftmargin=1.6em, itemsep=0.2em]
\item If $H\downarrow S_0$ (so $a\downarrow0$), then $V_0\to e^{-rT}$ (almost sure hit at time 0).
\item If $H\uparrow\infty$ (so $a\to\infty$), then $V_0\to 0$.
\end{itemize}

% ------------------------------------------------------------
\item \textbf{(Optional) Monte Carlo verification (discrete monitoring).}
Simulate $S$ on a grid $0=t_0<\cdots<t_N=T$ and approximate
\[
\mathbf 1_{\{\sup_{0\le t\le T} S_t\ge H\}}
\approx
\mathbf 1_{\{\max_{0\le k\le N} S_{t_k}\ge H\}}.
\]
Compute the Monte Carlo estimator
\[
\widehat V_0^{(N)} := e^{-rT}\,\frac1M\sum_{m=1}^M
\mathbf 1_{\{\max_k S_{t_k}^{(m)}\ge H\}},
\]
and compare it to the closed-form value from Question 3.
Comment on the bias as $N$ increases (discrete monitoring misses barrier hits between grid points).

\end{enumerate}

\end{exercise}


% ============================================================
% Additional Exercises --- Static Arbitrage (Lecture 9)
% ============================================================

\section{Lecture 8\& 9 --- Static arbitrage and consistency checks}

\begin{exercise}[Call spread bounds]
Fix a maturity \(T\) and two strikes \(K_1<K_2\). Let
\[
C(K_1), \qquad C(K_2)
\]
be the prices of European call options.

\begin{enumerate}[leftmargin=1.6em]
\item Write explicitly the payoff of the call spread
\[
(S_T-K_1)^+-(S_T-K_2)^+.
\]

\item Show that for all \(S_T\),
\[
0 \le (S_T-K_1)^+-(S_T-K_2)^+ \le K_2-K_1.
\]

\item Deduce that
\[
0 \le C(K_1)-C(K_2) \le e^{-r(T-t)}(K_2-K_1).
\]

\item Construct an arbitrage if one of these inequalities is violated.
\end{enumerate}
\end{exercise}

% ------------------------------------------------------------

\begin{exercise}[Monotonicity and arbitrage]
Let \(K_1<K_2\) and suppose that
\[
C(K_1) < C(K_2).
\]

\begin{enumerate}[leftmargin=1.6em]
\item Construct a portfolio involving two calls that yields a non-negative payoff.

\item Compute its initial cost.

\item Show that this leads to an arbitrage opportunity.
\end{enumerate}
\end{exercise}

% ------------------------------------------------------------

\begin{exercise}[Convexity and butterfly arbitrage]
Let \(K_1<K_2<K_3\) and suppose that
\[
C(K_2) > \frac{K_3-K_2}{K_3-K_1}C(K_1)
+
\frac{K_2-K_1}{K_3-K_1}C(K_3).
\]

\begin{enumerate}[leftmargin=1.6em]
\item Consider the butterfly portfolio
\[
(S_T-K_1)^+ - 2(S_T-K_2)^+ + (S_T-K_3)^+.
\]

\item Show that its payoff is non-negative.

\item Compute its initial cost under the above assumption.

\item Conclude that this leads to an arbitrage.
\end{enumerate}
\end{exercise}

% ------------------------------------------------------------

\begin{exercise}[Put--call parity violation]
Let \(C_t\) and \(P_t\) be the prices of a European call and put with the same strike \(K\) and maturity \(T\).

Suppose that
\[
C_t - P_t > S_t - K e^{-r(T-t)}.
\]

\begin{enumerate}[leftmargin=1.6em]
\item Identify which portfolio is overpriced.

\item Construct explicitly a trading strategy that yields an arbitrage.

\item Describe both the initial cashflow and the terminal payoff.
\end{enumerate}
\end{exercise}

% ------------------------------------------------------------

\begin{exercise}[Calendar arbitrage]
Consider two European calls with the same strike \(K\) but maturities \(T_1<T_2\).

Suppose that
\[
C(K,T_2) < C(K,T_1).
\]

\begin{enumerate}[leftmargin=1.6em]
\item Construct a trading strategy exploiting this violation.

\item Compute the initial cashflow.

\item Show that the value of the portfolio at time \(T_1\) is non-negative.

\item Conclude that this yields an arbitrage.
\end{enumerate}
\end{exercise}


% ============================================================
% Additional Exercises --- American Options
% ============================================================

\section{Lecture 10 --- American options}

% ------------------------------------------------------------

\begin{exercise}[American put in a two-step binomial model]
Consider a two-period binomial model with
\[
S_0=100,\qquad u=1.2,\qquad d=0.8,\qquad r=0,
\]
and an American put with strike \(K=100\).

\begin{enumerate}[leftmargin=1.6em]
\item Compute the risk-neutral probability \(p\).

\item Build the stock-price tree.

\item Compute the terminal payoff of the American put at maturity.

\item Using backward induction, compute the option value at each node.

\item At each node, compare the immediate exercise value with the continuation value and determine whether early exercise is optimal.
\end{enumerate}
\end{exercise}

% ------------------------------------------------------------

\begin{exercise}[American vs European put in a three-step tree]
Consider a three-period binomial model with
\[
S_0=100,\qquad u=1.1,\qquad d=0.9,\qquad 1+r=1.02,
\]
and strike \(K=100\).

\begin{enumerate}[leftmargin=1.6em]
\item Compute the price of the corresponding European put.

\item Compute the price of the corresponding American put.

\item Compare the two prices and compute the early exercise premium.

\item Identify the nodes where immediate exercise is optimal for the American put.
\end{enumerate}
\end{exercise}

% ------------------------------------------------------------

\begin{exercise}[Dynamic programming on a finite-state Markov chain]
Let \(X_n\) be a Markov chain on the state space
\[
\{1,2,3\}
\]
with transition matrix
\[
A=
\begin{pmatrix}
0.6 & 0.4 & 0 \\
0.2 & 0.5 & 0.3 \\
0 & 0.3 & 0.7
\end{pmatrix}.
\]
Assume the discount factor is \(\alpha=0.95\), maturity \(T=3\), and reward function
\[
g(1)=0,\qquad g(2)=2,\qquad g(3)=5.
\]

\begin{enumerate}[leftmargin=1.6em]
\item Compute \(V_3\).

\item Using the recursion
\[
V_n = \max\big(\alpha^n g(X_n),\, A V_{n+1}\big),
\]
compute \(V_2\), \(V_1\), and \(V_0\) componentwise.

\item For each state and each time, determine whether it is optimal to stop or continue.

\item Interpret economically the resulting exercise policy.
\end{enumerate}
\end{exercise}

% ------------------------------------------------------------

\begin{exercise}[Risk-neutral representation in continuous time]
Consider an American put in the Black--Scholes model with payoff
\[
g(s)=(K-s)^+.
\]

\begin{enumerate}[leftmargin=1.6em]
\item Write the continuous-time risk-neutral pricing representation
\[
V_t={\rm esssup}_{\tau\in\mathcal{T}_{t,T}}
\E^{\mathbb{Q}}\!\left[e^{-r(\tau-t)}(K-S_\tau)^+\mid\mathcal{F}_t\right].
\]

\item Write the corresponding European pricing formula.

\item Explain, using these two formulas, why
\[
V_t^{\mathrm{Am}}\ge V_t^{\mathrm{Eur}}.
\]

\item For the numerical values
\[
S_t=95,\qquad K=100,\qquad r=0.05,\qquad T-t=1,
\]
compute the immediate exercise value.

\item Explain qualitatively when early exercise may be optimal.
\end{enumerate}
\end{exercise}

% ------------------------------------------------------------

\begin{exercise}[Variational inequality for the American put]
Let \(v(t,s)\) denote the price of an American put in the Black--Scholes model.

\begin{enumerate}[leftmargin=1.6em]
\item Write the variational inequality satisfied by \(v\):
\[
\max\!\left((K-s)^+-v,\; v_t + rs\,v_s + \frac12 \sigma^2 s^2 v_{ss} - rv\right)=0.
\]

\item Specialize this formula to the parameters
\[
K=100,\qquad r=0.02,\qquad \sigma=0.2.
\]

\item Compute explicitly the terminal condition \(v(T,s)\) at
\[
s=60,\ 80,\ 100,\ 120.
\]

\item For each of these values of \(s\), compute the immediate exercise payoff.

\item Explain which of these values are more likely to lie in the stopping region and which are more likely to lie in the continuation region.
\end{enumerate}
\end{exercise}

% ------------------------------------------------------------

\begin{exercise}[Finite-difference / Longstaff--Schwartz intuition]
This exercise is computational and qualitative.

\begin{enumerate}[leftmargin=1.6em]
\item In a finite-difference method for an American put, explain what is done at each time step after computing the PDE update:
\[
v^{\text{new}}(t,s)
\quad \longrightarrow \quad
\max\big(v^{\text{new}}(t,s),(K-s)^+\big).
\]

\item Suppose at some grid point \((t,s)\) one obtains numerically
\[
v^{\text{new}}(t,s)=7.4
\qquad\text{and}\qquad
(K-s)^+=9.
\]
What value should be kept for the American option?

\item In the Longstaff--Schwartz method, one estimates the continuation value by regression.
Suppose that at one date the estimated continuation value is \(5.2\) and the immediate exercise payoff is \(6.1\).
Should one exercise or continue?

\item Suppose instead that the continuation value is \(8.4\) and the immediate payoff is \(6.1\).
What is the optimal decision?

\item Explain in one sentence the common idea behind the tree method, the finite-difference method, and the Longstaff--Schwartz method.
\end{enumerate}
\end{exercise}

% \section{Lecture 7 --- Risk Neutral Pricing and the Black Scholes Model.}


% % --- Put ~10--15 exercises here ---

% \section{Lecture 7 --- Connection between PDEs and Martingale representation.}


% % ==========================================================
% \section*{Appendix (optional)}
% % You can add formula sheets, common identities, etc.

\end{document}